\subsection{Widlar Mirror}\label{sec:mirror_widlar}

\subsubsection{Overview}
Constructed from two transistors and a source (emitter) degeneration
resistor \(G_S\) under the mirroring transistor. The resistor eats part
of \(V_{GS}\), so a small output current can be produced from a large
input current without an extreme \(\frac{W}{L}\) ratio, and the local
series feedback raises the output impedance.

\subsubsection{Schematic}
\begin{circuitfig}
  \draw (0,0) node[nmos, xscale=-1](M0){};
  \node[left] at (-0.4,0) {$M_0$};
  \draw (3,0) node[nmos](M1){};
  \node[right] at (3.4,0) {$M_1$};
  \draw (M0.G) -- (M1.G);
  \draw (M0.D) -- (0,1.5) node[circ]{} node[left]{$V_x$} -| (1.5,0) node[circ]{};
  \draw (0,1.5) to[european resistor, l={$G_2$}] (0,3.5) node[vcc]{$V_{cc}$};
  \draw (M1.D) -- (3,1.5) node[circ]{} node[right]{$V_{out}$} to[european resistor, l_={$G_3$}] (3,3.5) node[vcc]{$V_{cc}$};
  \draw (M1.S) -- (3,-1.2) node[circ]{} node[right]{$V_s$} to[european resistor, l_={$G_S$}] (3,-2.9);
  \draw (M0.S) -- (0,-2.9) -- (3,-2.9) -- (4,-2.9) node[right]{$V_{ee}$};
\end{circuitfig}
As in Section~\ref{sec:mirror_basic}, \(G_3\) exists only for analysis
(the load).

\subsubsection{Transfer Function}
\begin{tabularx}{\linewidth}{|X|X|X|X|X|X|} \hline
  Transistor & Gate (G) & Source (S) & Drain (D) & $V_{GS}$ & $V_{DS}$ \\ \hline
  $M_0$ & $V_{x}$ & $V_{ee}$ & $V_{x}$ & $V_{x} - V_{ee}$ & $V_{x} - V_{ee}$ \\ \hline
  $M_1$ & $V_{x}$ & $V_{s}$ & $V_{out}$ & $V_{x} - V_{s}$ & $V_{out} - V_{s}$ \\ \hline
\end{tabularx}
\begin{enumerate}
  \item \textbf{Large Signal, Time Domain, DC Input Analysis (Q5)}: (\(K_{CLM} = 0\))
    \begin{align*}
      I_{in} = I_{M_0} &= k_{n_0}(V_x - V_{ee} - V_{th})^2
        \hspace{0.5cm} \Rightarrow \hspace{0.5cm} V_x - V_{ee} = V_{th} + \sqrt{\frac{I_{in}}{k_{n_0}}} \tag{1} \\
      I_{out} = I_{M_1} &= k_{n_1}(V_x - V_s - V_{th})^2
        \hspace{0.5cm} \Rightarrow \hspace{0.5cm} V_x - V_s = V_{th} + \sqrt{\frac{I_{out}}{k_{n_1}}} \tag{2} \\
      (V_s - V_{ee}) G_S &= I_{out}
        \hspace{0.5cm} \Rightarrow \hspace{0.5cm} V_s - V_{ee} = \frac{I_{out}}{G_S} \tag{3}
    \end{align*}
    Each square law is solved for its gate-source voltage, e.g.\ for (1):
    \begin{align*}
      (V_x - V_{ee} - V_{th})^2 &= \frac{I_{in}}{k_{n_0}} \\
      V_x - V_{ee} - V_{th} &= \sqrt{\frac{I_{in}}{k_{n_0}}} \quad \text{[positive root, } M_0 \text{ on]}
    \end{align*}
    Adding (2) and (3), the left sides sum to \((V_x - V_s) + (V_s - V_{ee}) = V_x - V_{ee}\),
    which is the left side of (1). Equating the right sides:
    \begin{align*}
      V_{th} + \sqrt{\frac{I_{out}}{k_{n_1}}} + \frac{I_{out}}{G_S} &= V_{th} + \sqrt{\frac{I_{in}}{k_{n_0}}} \\
      \frac{I_{out}}{G_S} + \sqrt{\frac{I_{out}}{k_{n_1}}} - \sqrt{\frac{I_{in}}{k_{n_0}}} &= 0 \tag{4}
    \end{align*}
    With \(u = \sqrt{I_{out}}\), so that \(I_{out} = u^2\) and
    \(\sqrt{\frac{I_{out}}{k_{n_1}}} = \frac{u}{\sqrt{k_{n_1}}}\), (4) is a quadratic
    \(a u^2 + b u + c = 0\) with
    \(a = \frac{1}{G_S}\), \(b = \frac{1}{\sqrt{k_{n_1}}}\), \(c = -\sqrt{\frac{I_{in}}{k_{n_0}}}\):
    \begin{align*}
      u &= \frac{-b + \sqrt{b^2 - 4ac}}{2a} \hspace{1cm} [u > 0 \text{, the other root is negative}] \\
        &= \frac{-\frac{1}{\sqrt{k_{n_1}}} + \sqrt{\frac{1}{k_{n_1}} + \frac{4}{G_S}\sqrt{\frac{I_{in}}{k_{n_0}}}}}{\frac{2}{G_S}} \\
        &= \frac{G_S}{2} \left( \sqrt{\frac{1}{k_{n_1}} + \frac{4}{G_S}\sqrt{\frac{I_{in}}{k_{n_0}}}} - \frac{1}{\sqrt{k_{n_1}}} \right) \\
      I_{out} = u^2 &= \frac{G_S^2}{4} \left( \sqrt{\frac{1}{k_{n_1}} + \frac{4}{G_S}\sqrt{\frac{I_{in}}{k_{n_0}}}} - \frac{1}{\sqrt{k_{n_1}}} \right)^2
    \end{align*}
    For design it is easier to invert (4) for \(G_S\). From the square
    law, \(\sqrt{\frac{I_{D}}{k_{n}}} = V_{GS} - V_{th} = V_{ov}\):
    \begin{align*}
      \frac{I_{out}}{G_S} &= \sqrt{\frac{I_{in}}{k_{n_0}}} - \sqrt{\frac{I_{out}}{k_{n_1}}} \\
      \frac{1}{G_S} = R_S &= \frac{\sqrt{\frac{I_{in}}{k_{n_0}}} - \sqrt{\frac{I_{out}}{k_{n_1}}}}{I_{out}} = \frac{V_{ov_0} - V_{ov_1}}{I_{out}}
    \end{align*}
    \(G_S \to \infty\) sends \(\frac{I_{out}}{G_S} \to 0\) in (4) and gives back the basic mirror:
    \begin{align*}
      \sqrt{\frac{I_{out}}{k_{n_1}}} &= \sqrt{\frac{I_{in}}{k_{n_0}}}
      \hspace{0.5cm} \Rightarrow \hspace{0.5cm} \frac{I_{out}}{I_{in}} = \frac{k_{n_1}}{k_{n_0}}
    \end{align*}
  \item \textbf{Small Signal, Frequency Domain, AC Input Analysis (Q4)}: (\(g_{mb}\) neglected)\\
    Applying KCL at $V_{x}$ with the input current \(I_{in}\) injected
    (\(M_0\) has \(V_{gs_0} = V_{ds_0} = V_x\), \(I_G = 0\)):
    \begin{align*}
      V_x G_2 + g_{m_0} V_x + g_{o_0} V_x - I_{in} &= 0 \\
      V_x(G_2 + g_{o_0} + g_{m_0}) &= I_{in} \tag{5}
    \end{align*}
    Applying KCL at $V_{s}$ (current into \(G_S\) equals drain current of \(M_1\)):
    \begin{align*}
      V_s G_S &= g_{m_1}(V_x - V_s) + g_{o_1}(V_{out} - V_s) \\
      V_s (G_S + g_{m_1} + g_{o_1}) &= g_{m_1} V_x + g_{o_1} V_{out} \tag{6}
    \end{align*}
    Applying KCL at $V_{out}$, and replacing the drain current of \(M_1\),
    \(g_{m_1}(V_x - V_s) + g_{o_1}(V_{out} - V_s)\), by \(V_s G_S\) (first line of (6)):
    \begin{align*}
      V_{out} G_3 + g_{m_1}(V_x - V_s) + g_{o_1}(V_{out} - V_s) &= 0 \\
      V_{out} G_3 + V_s G_S &= 0 \tag{7}
    \end{align*}
\end{enumerate}

\subsubsection{Analysis}
\begin{enumerate}
  \item \textbf{Current Gain, Small Signal, Frequency Domain, AC Input (Q4)}:
    output at AC ground (\(V_{out} = 0\)); the drain current of \(M_1\) is
    the current in \(G_S\), \(I_{out} = V_s G_S\). From (6) with \(V_{out} = 0\) and (5):
    \begin{align*}
      V_s (G_S + g_{m_1} + g_{o_1}) &= g_{m_1} V_x + g_{o_1} \cdot 0 \\
      V_s &= \frac{g_{m_1} V_x}{G_S + g_{m_1} + g_{o_1}} \\
      V_x &= \frac{I_{in}}{G_2 + g_{o_0} + g_{m_0}} \quad \text{[from (5)]} \\
      I_{out} = V_s G_S &= \frac{G_S\, g_{m_1}}{G_S + g_{m_1} + g_{o_1}} \cdot \frac{I_{in}}{G_2 + g_{o_0} + g_{m_0}} \\
      \frac{I_{out}}{I_{in}} &= \frac{g_{m_1}}{1 + (g_{m_1} + g_{o_1}) R_S} \cdot \frac{1}{G_2 + g_{o_0} + g_{m_0}} \quad \text{[dividing by } G_S = \tfrac{1}{R_S} \text{]} \\
                             &\approx \frac{g_{m_1}}{g_{m_0}} \cdot \frac{1}{1 + g_{m_1} R_S} \hspace{1cm} [g_{o_1} \ll g_{m_1},\ G_2, g_{o_0} \ll g_{m_0}]
    \end{align*}
  \item \textbf{Impedances}:
    \paragraph{At Input:} identical to the basic mirror, from (5):
      \[
        Z_{in} = \frac{1}{G_2 + g_{o_0} + g_{m_0}} \approx \frac{1}{g_{m_0}}
      \]
    \paragraph{At Output:} \(I_{in} = 0 \Rightarrow V_x = 0\) by (5). Apply
      \(I_{test}\) at \(V_{out}\) (with \(G_3\) removed); it flows through
      \(M_1\) into \(G_S\): \(V_s = \frac{I_{test}}{G_S}\). From (6):
      \begin{align*}
        \frac{I_{test}}{G_S}(G_S + g_{m_1} + g_{o_1}) &= g_{o_1} V_{out} \\
        I_{test}\left(1 + \frac{g_{m_1} + g_{o_1}}{G_S}\right) &= g_{o_1} V_{out} \\
        Z_{out} = \frac{V_{out}}{I_{test}} &= \frac{1 + (g_{m_1} + g_{o_1}) R_S}{g_{o_1}} \\
                &= \frac{1}{g_{o_1}} + \frac{g_{m_1} R_S}{g_{o_1}} + \frac{g_{o_1} R_S}{g_{o_1}} \\
        Z_{out} &= r_{o_1} + g_{m_1} r_{o_1} R_S + R_S \approx r_{o_1}(1 + g_{m_1} R_S) \quad \text{[} R_S \ll r_{o_1} \text{]}
      \end{align*}
      The degeneration multiplies the basic mirror output impedance by
      \(1 + g_{m_1} R_S\).
  \item \textbf{Output Load (Max), Large Signal, Time Domain, DC Input (Q5)}:
    \(M_1\) saturated needs \(V_{DS_1} = V_{out} - V_s \geq V_{ov_1}\), with
    \(V_s = V_{ee} + \frac{I_{out}}{G_S} = V_{ee} + I_{out} R_S\) from (3):
    \begin{align*}
      V_{out} &\geq V_s + V_{ov_1} \\
      V_{out,\min} &= V_{ee} + I_{out} R_S + V_{ov_1} \\
      R_{L,\max} &= \frac{V_{cc} - V_{out,\min}}{I_{out}}
    \end{align*}
\end{enumerate}

\subsubsection{Design}
\begin{enumerate}
  \item Choose \(V_{ov_0}\); then \(V_x = V_{ee} + V_{th} + V_{ov_0}\) and
    \(W_0 = \frac{2 I_{in} L_0}{K_P V_{ov_0}^2}\).
  \item Choose \(W_1\) (often \(W_1 = W_0\)); then from
    \(I_{out} = \frac{K_P}{2}\frac{W_1}{L_1} V_{ov_1}^2\),
    \(V_{ov_1} = \sqrt{\frac{2 I_{out} L_1}{K_P W_1}}\).
  \item \(R_S = \frac{V_{ov_0} - V_{ov_1}}{I_{out}}\) and
    \(R_2 = \frac{V_{cc} - V_x}{I_{in}}\).
\end{enumerate}

\subsubsection{Question}
Generate \(I_{out} = 10\,\mu\)A from \(I_{in} = 100\,\mu\)A with
\(W_0 = W_1\), \(L = 1\,\mu\)m, \(K_P = 200\,\mu\)A/V\(^2\), \(V_{th} = 1\)V,
\(K_{CLM} = 0.02\)V\(^{-1}\), \(V_{cc} = 5\)V, \(V_{ee} = 0\)V,
\(V_{ov_0} = 0.2\)V. Find \(R_S\), \(R_2\), \(Z_{out}\) and compare with a
basic mirror delivering the same \(10\,\mu\)A.
\paragraph{Solution:}
\begin{align*}
  W_0 &= \frac{2 \cdot 100\,\mu A \cdot 1\,\mu m}{200\,\mu A/V^2 \cdot (0.2V)^2} = 25\,\mu m = W_1 \\
  k_{n_1} &= \frac{K_P}{2}\frac{W_1}{L_1} = 100\,\mu A/V^2 \cdot 25 = 2.5\,mA/V^2 \\
  V_{ov_1} &= \sqrt{\frac{I_{out}}{k_{n_1}}} = \sqrt{\frac{10\,\mu A}{2.5\,mA/V^2}} = \sqrt{0.004\,V^2} = 63.2\,mV \\
  R_S &= \frac{V_{ov_0} - V_{ov_1}}{I_{out}} = \frac{0.2V - 0.0632V}{10\,\mu A} = \frac{0.1368V}{10\,\mu A} = 13.68\,k\Omega \\
  V_x &= V_{ee} + V_{th} + V_{ov_0} = (0 + 1 + 0.2)V = 1.2V \\
  R_2 &= \frac{V_{cc} - V_x}{I_{in}} = \frac{5V - 1.2V}{100\,\mu A} = 38\,k\Omega \\
  g_{m_1} &= \frac{2 I_{out}}{V_{ov_1}} = \frac{2 \cdot 10\,\mu A}{0.0632V} = 316\,\mu S \hspace{1cm}
  r_{o_1} = \frac{1}{K_{CLM} I_{out}} = \frac{1}{0.02 \cdot 10\,\mu A} = 5\,M\Omega \\
  g_{m_1} R_S &= 316\,\mu S \cdot 13.68\,k\Omega = 4.32 \\
  Z_{out} &= r_{o_1} + R_S + g_{m_1} r_{o_1} R_S = 5\,M\Omega + 13.68\,k\Omega + 4.32 \cdot 5\,M\Omega \\
          &= 5\,M\Omega + 0.014\,M\Omega + 21.6\,M\Omega \approx 26.6\,M\Omega \\
  Z_{out,basic} &= r_{o_1} = 5\,M\Omega \hspace{1cm} \frac{Z_{out}}{Z_{out,basic}} \approx 1 + g_{m_1} R_S = 5.32 \\
  V_{out,\min} &= V_{ee} + I_{out} R_S + V_{ov_1} = 0 + 10\,\mu A \cdot 13.68\,k\Omega + 0.0632V = 0.1368V + 0.0632V = 0.2V \\
  R_{L,\max} &= \frac{V_{cc} - V_{out,\min}}{I_{out}} = \frac{5V - 0.2V}{10\,\mu A} = 480\,k\Omega
\end{align*}
\(V_{ov_1} = 63\,\)mV is close to weak inversion where the square law
overestimates \(I_D\); the result should be checked in Q5 simulation
(DC/OP).
